\section{A small power key grid}
\label{sec:a-small-power-key-grid}

We next exclude nearly pure patches whose two laws both have small power
width. The construction uses a separate anchor for each cell in a grid.
We first arrange that an anchor with a sufficiently spread-out law is
unlikely to have very small degree into a narrow law.

Take rational constants \(0<D_0<\min(.1,D_*/2)\), where \(D_*\) is from
Lemma~\ref{lem:small-polynomial-broad-side}; here \(D_0\) is an exponent,
not a height step length. Fix
\(0<d<D_0/1000,\ b=d/4,\ c=b/20\). We claim that, after stabilization
and \(o(N)\) discards, we may assume
\begin{equation}
 \begin{gathered}
 \text{every }\sigma\text{ on }X,\ \tau\text{ on }Y\text{ with }\quad
 N\max\sigma\le n^{D_0},\quad \operatorname{width}(\tau)\le n^{1/4}\\
 \text{have both colors of density }> n^{-c}.
 \end{gathered}\label{eq:source-1}
\end{equation}
To prove the claim, first suppose that a family violating this assertion
were available. On an available subfamily one color \(G\) has density at
least \(1-n^{-c}\). For each witness put
\[
 B_\sigma=\{x:d_G(x,\tau)<1-n^{-c/2}\},\qquad
 \sigma'=\sigma|_{X\setminus B_\sigma}.
\]
Markov's inequality gives \(\sigma(B_\sigma)\le n^{-c/2}\), so this
restriction increases the law by at most a factor of 2. In particular,
\(N\max\sigma'\le2n^{D_0}\le n^{D_*}\) for large \(n\), and every
label in its support has degree at least \(1-n^{-c/2}\) into \(\tau\).
The restricted witnesses are still available. Choose a balanced mixture
by Lemma~\ref{lem:balanced-profiles}, and apply
Lemma~\ref{lem:small-polynomial-broad-side} with the sides interchanged:
its narrow law is \(\tau\), its broad law is \(\sigma'\), and its defect
exponent is \(p_0=c/2\). This gives a monochromatic cube, a
contradiction. Thus violations are not available, and
Lemma~\ref{lem:stabilization} makes them absent after the claimed
subsequence and discards.

\begin{lemma}[Small-grid purity exclusion]\label{lem:small-grid-purity}
With the constants fixed above, on sets satisfying \eqref{eq:source-1}
there cannot be an available family of one-color patches
\((\mu_i,\nu_i)\) of widths \(\le n^{d/2}\) on \(X,Y\), with every
\(y\in\operatorname{supp}\nu_i\) having degree into \(\mu_i\) in the
given color at least \(1-e^{-n^{p}}\) for fixed \(p>0\).
\end{lemma}

\begin{proof}
\noindent\textbf{Step 1: the grid and its raw anchor laws.}\quad
Suppose such a family exists. A finite menu retains availability by
choosing a representative for each successful discard pair. Denote its
given color by \(G\).

Partition some cube coordinates into \(s=\lceil n^d\rceil\) chunks of
\(\lfloor n^d\rfloor\) special bits each, with total length
\(m=\Theta(n^{2d})\), and set aside \(q=\lceil n^{4d}\rceil\) further
auxiliary bits. Bin each special chunk's count into consecutive intervals
of probability at most \(2n^{-b}\); this is possible since each binomial
atom is \(O(n^{-d/2})\). The vector of bin indices is a grid key \(g\),
and the auxiliary word is a second key \(t\in Q_q\). Write \(E(g)\) for
the distinct grid keys obtained by moving one bin in one coordinate, so
\(|E(g)|\le2s\). Grid distance means the number of such one-bin moves.
An actual cube position reads the cell indexed by its
pair \((g,t)\); many positions can read the same cell. A fixed \(g\)
occurs with fraction at most \((2n^{-b})^s\). Since \(m+q=o(n)\), unused
residual bits remain, and \(t\) is uniform on either parity of the cube,
also when \(g\) is fixed.

For each grid key independently sample a patch tag \(i_g\). Conditional
on these tags, sample independent anchors
\(W_{g,t}\sim\mu_{i_g}\), one for every cell. The distributions of the
tags will be chosen below. For odd positions of key \((g,t)\), restrict
\(\nu_{i_g}\) to labels adjacent to all anchors in the two lists
\[
 W_{h,t}\quad(h\in E(g)),\qquad
 W_{g,t'}\quad(d_H(t,t')\le1).
\]
Call these the cross list and the own list, respectively. They cover all
incident even anchors: a special-bit flip either leaves \(g\) unchanged
or moves it to \(E(g)\), an auxiliary-bit flip changes only \(t\), and a
residual-bit flip changes neither key.

We need to be able to omit any specified cross hit later. For this
reason, process the cross list once with each possible element last,
using a fixed order of the other elements independent of \(t\). In every
such order require each successive restriction to retain a fraction at
least \(n^{-c}\). The empty cross list is allowed. After the full cross
restriction, require the own list, taken together, to retain a fraction
at least \(1-n^{-2}\). Call the cell valid when all these requirements
hold. On validity let \(p_{g,t}\) be the normalized fully restricted
law; otherwise let it be zero. Thus it is a subprobability supported on
hits of all incident even anchors, and
\[
 Np_{g,t}\le L:=\exp(n^{d/2}+2c s\log n+O(1)).
\]

For an anchor variable \(w\) in the lists, let \(p_{g,t}^{(-w)}\) be
\(\nu_{i_g}\) restricted to all listed hits except the hit of \(w\),
then normalized. If the restriction has zero mass, use a fixed fallback
independent of \(w\). Deletion is by the variable's name, even if two
anchors happen to take the same label. On validity,
\[
 \begin{aligned}
 p_{g,t}&\le(1-n^{-2})^{-1}p_{g,t}^{(-w)}
       &&\text{for an own hit},\\
 p_{g,t}&\le n^c(1-n^{-2})^{-1}p_{g,t}^{(-w)}
       \le n^{2c}p_{g,t}^{(-w)}
       &&\text{for a cross hit}.
 \end{aligned}
\]
The second comparison uses the order in which that cross hit is last,
followed by the restriction to the own list. Each deletion law is
normalized without consulting the deleted variable's value.

Choose product profiles for the tags so that, for a constant \(K'\),
\begin{enumerate}[label=\textup{(\roman*)}]
\item \(\mathbb E N\mu_{i_g}\le K'\) pointwise for every grid key \(g\);
\item for each \(g\), the expectation of \(Np_{g,t}\), averaged over
odd sites with this grid key, is pointwise \(\le K'\).
\end{enumerate}
Both expectations refer to the raw independent tag and anchor model.
To obtain these profiles, use one player per grid key in
Lemma~\ref{lem:balanced-profiles}. Fix the other players' profiles and
assign nonnegative prices to the two output vectors of this player.
Availability permits a patch avoiding, on each side, all labels whose
price exceeds a sufficiently large constant times total price divided
by \(N\). Its first output \(\mu_{i_g}\) is supported on its first
support; every possible second output \(p_{g,t}\) is a subprobability on
its second support. Both price contributions therefore have the required
bound. Separation gives a mixture satisfying the two coordinatewise
bounds, and the simultaneous-profile part of the lemma gives profiles
for all players. Its hypothesis on the output rule holds because the
filter and validity rules do not use the profile probabilities.

\noindent\textbf{Step 2: filter tails and tag conditioning.}\quad
Consider one step in one prescribed cross order, conditional on the
starting tag and the preceding cross anchors, all of whose tests have
passed. The current filtered law \(\tau\) has width at most
\(\log L<n^{1/4}\). The next anchor comes from a distinct grid key, so
its raw marginal law \(\alpha\) is independent of this history and
satisfies \(N\alpha\le K'\) pointwise. Put
\[
 H_\tau=\{x:d_G(x,\tau)<n^{-c}\}.
\]
If \(\alpha(H_\tau)\ge K'n^{-D_0}\), the normalized restriction
\(\sigma=\alpha|_{H_\tau}\) satisfies
\(N\max\sigma\le n^{D_0}\), but
\(d_G(\sigma,\tau)<n^{-c}\), contrary to \eqref{eq:source-1}. Thus each
cross step fails with probability \(O(n^{-D_0})\). There are at most
\(O(s^2)\) order--step pairs, giving total cross-failure probability
\[
 O(s^2n^{-D_0})=n^{-D_0+O(d)+o(1)}.
\]

For the own test, condition on all tags and the cross anchors, with the
cross tests passed. Each label in the support of the resulting law has
degree at least \(1-e^{-n^p}\) into \(\mu_{i_g}\). The own anchors are
still independent draws from \(\mu_{i_g}\). The expected mass removed by
their \(q+1\) constraints is at most \((q+1)e^{-n^p}\), so Markov's
inequality bounds own-test failure by \(n^2(q+1)e^{-n^p}\).

We first impose a condition on the tags alone. Let \(B_g\) be the event
that, conditional on all tags, the cross-failure probability for cell
\((g,t)\), including all orders and steps, exceeds \(n^{-D_0/2}\).
This probability is independent of \(t\) by auxiliary-key symmetry.
The event \(B_g\) reads just the tags at \(g\) and \(E(g)\), and
\[
 \Pr(B_g)\le O(s^2n^{-D_0})n^{D_0/2}
 =n^{-D_0/2+2d+o(1)}\le n^{-D_0/3}.
\]
The overlap graph of these tag events has degree \(O(s^2)\).
Lemma~\ref{lem:conditional-avoidance} therefore gives a tag law avoiding
all \(B_g\), with charges \(n^{-D_0/4}\), since
\(O(s^2)n^{-D_0/4}=o(1)\).

This tag law also has bounded average anchor loads with probability
\(1-o(1)\). Here is the moment calculation. For a fixed \(x\in X\),
use even roles \(v\in A\) and weights
\(Z_v=N\mu_{i_{g(v)}}(x)\), of cap \(L_\mu=e^{n^{d/2}}\). Declare
roles near when their grid keys are within a fixed sufficiently large
distance. Each near set has fraction
\(f\le n^{O(1)}(2n^{-b})^s\), and \(nfL_\mu=o(1)\). For \(l\le n\)
separated roles, remove the tag-avoidance constraints touching their tag
variables. There are \(O(ls)\) such constraints, so the comparison cost
is
\[
 \exp\big(O(lsn^{-D_0/4})\big)=e^{o(l)}.
\]
The remaining integral uses independent raw tags, each with mean weight
at most \(K'\) by (i). Lemma~\ref{lem:scattered-moments}, followed by
Markov's inequality and a union over \(x\), bounds
\(|A|^{-1}\sum_{v\in A}N\mu_{i_{g(v)}}(x)\) by a constant
simultaneously for all \(x\). Call such tags typical.

\noindent\textbf{Step 3: predictive alarms before anchor conditioning.}\quad
We next ensure that the eventual odd labels give a sufficiently large
predictive denominator at every even role. Fix an even position \(v\)
of key \((g,t)\), all tags, and all anchors except
\(z=W_{g,t}\). For a vector \(\mathbf y=(y_u:u\sim v)\) of odd labels
put
\[
 F_z(\mathbf y)=1_{\mathcal V}\prod_{u\sim v}p_{g(u),t(u)}(y_u),\qquad
 M_v(\mathbf y)=\int F_z(\mathbf y)\,d\mu_{i_g}(z),
\]
where \(\mathcal V\) requires validity of all the neighboring filters.
Both validity and the kernels are recomputed as \(z\) varies. Let
\(Q_v\) be the product of their deletion laws for the named variable
\(W_{g,t}\), with the fallbacks above when necessary. This is a
probability law on \(\mathbf y\) independent of \(z\). At most \(m\)
actual neighbors impose a cross hit of \(z\), so the deletion bounds give
\[
 F_z\le\exp\big(2cm\log n+O(n^{-1})\big)Q_v=e^{o(q)}Q_v;
 \qquad m\log n=o(q).
\]
Declare predictive failure when \(M_v=0\) or
\(M_v<e^{-.04q}Q_v\).

For fixed realized anchors \(\mathbf W\), let \(r_v(\mathbf W)\) be
the probability of this failure under independent neighbor draws,
multiplied by \(1_{\mathcal V}\); set it to zero if a neighboring filter
is invalid. Invalid rows may use arbitrary local fallback probabilities
in this auxiliary draw. Before any anchor conditioning, the data
subdensity on \(\mathcal V\) is exactly \(M_v\). Hence, even with the
tags and all other anchors fixed,
\[
 \mathbb E_{W_{g,t}}r_v\le e^{-.04q},\qquad
 \Pr_{W_{g,t}}(r_v>e^{-.02q})\le e^{-.02q}.
\]
These estimates retain \(\mathcal V\) as an indicator; they do not
condition the raw law on validity.

At a given key \((g,t)\) we must impose the inequality
\(r_v\le e^{-.02q}\) for all neighbor multiplicities that can occur at
an even position with that key. Each auxiliary flip contributes once.
The special-bit neighbors are described by at most \(2s+1\) counts:
one for each possible adjacent grid bin and one for the unchanged key.
The number of residual-bit neighbors is fixed. Thus there are at most
\((n+1)^{2s+1}=\exp(O(s\log n))\) distinct formulas, up to reordering
the neighbor labels. A union over these formulas, rather than over cube
positions, bounds the probability that a key violates an alarm by
\[
 \exp\big(-.02q+O(s\log n)\big)=\exp(-\Omega(q)),
\]
uniformly in the tags.

\noindent\textbf{Step 4: anchor conditioning and odd loads.}\quad
Fix tags admitted by Step 2. Group, for each \((g,t)\), its filter
failure and all its predictive alarms into one bad event on the raw
independent anchors. Give the grid-times-auxiliary graph edges for one
grid-bin move or one auxiliary-bit flip. Each grouped event reads
anchors within distance two of its key: the filter reads distance one,
and an alarm reads the filters at neighboring keys. The event graph
therefore has degree \(O((s+q)^4)\). Its failure probability is at most
\[
 n^{-D_0/2}+n^2(q+1)e^{-n^p}+\exp(-\Omega(q)).
\]
Charges \(n^{-D_0/4}\) again satisfy
Lemma~\ref{lem:conditional-avoidance}, since \(16d<D_0/4\). Use the
resulting anchor law at each fixed admitted tag history. Sampling in
these two stages leaves the tag law unchanged; the anchor avoidance
probability is not used to reweight it.

We show that the odd kernels have small column sums under this law. Fix
\(y\in Y\), and for an odd role \(u\) put
\[
 Z_u=Np_{g(u),t(u)}(y),\qquad
 a_u(\mathbf i)=\mathbb E_{\mathrm{raw\ anchors}\mid\mathbf i}Z_u,
 \qquad d_u=\mathbb E_{\mathrm{raw\ tags}}a_u(\mathbf i).
\]
The cap is \(L\). Take the near relation to mean that the grid keys are
within a fixed distance large enough to contain the local scopes just
described. Its fraction is again
\(f\le n^{O(1)}(2n^{-b})^s\), and
\[
 nfL\le\exp\big(-(b-2c+o(1))s\log n\big)=o(1),
\]
because \(b>2c\).

For \(l\le n\) separated roles first integrate the anchors, keeping the
tags fixed. To recover the independent raw anchor means, remove all
avoidance constraints touching their anchor lists. A list has radius
one and a constraint has radius two, so at most \(O((s+q)^3)\)
constraints touch each list. The lists are disjoint, and the resulting
integral factors as \(\prod_j a_{u_j}(\mathbf i)\). Each factor uses
only tags within grid distance one of its role. Next remove the
tag-avoidance constraints touching these tag lists, at most \(O(s^2)\)
per role, and integrate the independent raw tags. Consequently the full
two-stage law satisfies
\[
 \mathbb E\prod_{j=1}^l Z_{u_j}
 \le \exp\!\left(O\!\left(l\big((s+q)^3+s^2\big)n^{-D_0/4}\right)\right)
       \prod_{j=1}^l d_{u_j}
 \le C^l\prod_{j=1}^l d_{u_j},
\]
for a constant \(C\) independent of \(n\). The exponent is \(o(l)\), since
\(12d<D_0/4\). Profile constraint (ii) gives
\(|B|^{-1}\sum_{u\in B}d_u\le K'\).

Lemma~\ref{lem:scattered-moments} now bounds the normalized average
\(|B|^{-1}\sum_u Z_u\) by a constant, simultaneously for all labels,
with probability \(1-o(1)\). Here and in Step 2 the union is legitimate
because \(N\le n2^n\): a sufficiently large fixed threshold in the
\(n\)-th moment bound beats all \(N\) columns. The actual column sum is
\[
 \sum_{u\in B}p_{g(u),t(u)}(y)
 =\frac{|B|}{N}\left(\frac1{|B|}\sum_{u\in B}Z_u\right)=o(1),
\]
since \(C_n=N/2^n\to\infty\). It is therefore at most the clock-lemma
constant \(\theta_0\) for large \(n\). This property and typicality of
the tags hold simultaneously with probability \(1-o(1)\).

\Needspace{3\baselineskip}
\noindent\textbf{Step 5: injective odd labels and posterior even rows.}\quad
On a prehistory with the odd column bound, apply
Lemma~\ref{lem:clock-sampling}, for example with \(B=3\), to the actual
odd-position kernels. Use predictive failure at each even position as
a forbidden predicate. Each predicate reads its \(n\) odd neighbors,
and each odd row belongs to \(n\) such scopes. Its product-law
probability is at most \(e^{-.02q}\) by the anchor alarms. The maximum
label atom is at most \(L/N\), with \(\log L=o(q)=o(n)\). These bounds
meet the lemma's polynomial requirements. We obtain an odd injection
avoiding all predictive failures, with joint upper comparison
\(1+o(1)\) on at most \(n^2\) queried rows.

For each even role \(v\), evaluated at its assigned neighbor data, use
the probability row
\[
 p_v^X(x)=\frac{F_x(\mathbf y)\mu_{i_g}(x)}{M_v(\mathbf y)},
 \qquad g=g(v).
\]
It is supported on common neighbors of those odd labels. Predictive
success and the bound for \(F_x/Q_v\) give
\(Np_v^X\le e^{.04q+o(q)}\), including the prior cap
\(N\mu_{i_g}\le e^{n^{d/2}}\). In auxiliary experiments set the row
to zero unless \(\mathcal V\) and predictive success hold.

\noindent\textbf{Step 6: posterior cancellation and Hall.}\quad
Fix typical tags and a label \(x\in X\). For the even load moments use
\(Z_v=Np_v^X(x)\), cap \(L_{\rm even}=e^{.04q+o(q)}\), and declare
roles near when their auxiliary words are within Hamming distance five.
Then
\[
 f\le n^{O(1)}2^{-q},\qquad
 nfL_{\rm even}\le n^{O(1)}2^{-q}e^{.04q+o(q)}=o(1).
\]
For \(l\le n\) separated roles their odd-neighbor sets are disjoint.
Keep the successful-prehistory indicator while applying the clock
lemma's joint comparison to these at most \(n^2\) odd labels. This
replaces their injection law by independent odd draws. We may now
remove nonlocal success restrictions, retaining each local gate
\(\mathcal V\) in its star likelihood.

Next remove the anchor-avoidance constraints touching the \(l\) target
variables \(W_{g(v),t(v)}\). There are at most \(O(l(s+q)^2)\) such
constraints, and their total comparison cost is
\(\exp(O(l(s+q)^2n^{-D_0/4}))=e^{o(l)}\), in particular \(O(1)^n\).
The calculation at any other retained role does not consult a target
being integrated: its anchor inputs have auxiliary words within distance
two of its own word, whereas retained words are separated by more than
five.

With all other anchors fixed, integrating one target and its product
neighbor data gives the marginal data subdensity \(M_v\). Thus, with
zero terms where \(M_v=0\),
\[
 \sum_{\mathbf y}M_v(\mathbf y)
       \frac{F_x(\mathbf y)\mu_{i_g}(x)}{M_v(\mathbf y)}
 =\mu_{i_g}(x)\sum_{\mathbf y}F_x(\mathbf y)
 \le\mu_{i_g}(x).
\]
The last inequality uses the local validity gate; retaining predictive
success can only decrease the integral. We can perform this cancellation
successively at all separated roles. If \(\mathcal S\) denotes the
successful prehistory and assignment, the resulting comparison is
\[
 \mathbb E\!\left[1_{\mathcal S}\prod_{j=1}^l Z_{v_j}
                 \,\middle|\,\mathbf i\right]
 \le C^l\prod_{j=1}^l N\mu_{i_{g(v_j)}}(x).
\]
The comparison means have bounded average because the tags are typical.
Lemma~\ref{lem:scattered-moments} and the same Markov and column union
therefore bound the probability of \(\mathcal S\) together with any
normalized even load exceeding a sufficiently large fixed constant by
\(o(1)\), uniformly over typical tags.
Average this estimate over the tag law. Step 4 gives typical tags and
the odd column bound with probability \(1-o(1)\), and the clock sampler
avoids all predictive failures on those prehistories. There is therefore
a realization with both odd and even column bounds. The even column
sums are \(O(2^n/N)=o(1)\), hence at most 1. Its rows are probability
laws, so Hall's theorem gives distinct even labels. Together with the
odd injection they form a monochromatic cube, proving the lemma.
\end{proof}

\begin{corollary}[Initial discrepancy]\label{cor:initial-discrepancy}
After passing to a subsequence of a bad sequence and discarding a total
of \(o(N)\) labels, there is a fixed \(\eta_0>0\) such that
for every pair of probability laws \(\mu,\nu\) on the retained sides
and either color \(G\),
\begin{equation}
 \operatorname{width}(\mu),\operatorname{width}(\nu)\le n^{\eta_0}
 \quad\Longrightarrow\quad
 |d_G(\mu,\nu)-1/2|\le n^{-\eta_0} \label{eq:source-2}
\end{equation}
\end{corollary}

\begin{proof}
Take \(\beta=\gamma=d/4\) in Lemma~\ref{lem:bias-purity}, with its
corresponding \(h>0\). Stabilize its bias property at widths
\(n^\gamma\) and its purity property at doubled width budgets. If purity were available at widths
\(2n^\gamma\) and defect \(\le e^{-n^h}\), first retain an available
color on a subsequence. Restrict its second law to columns of defect at
most \(e^{-n^{h/2}}\); Markov's inequality loses only \(o(1)\) mass.
The restricted family is still available and has both widths at most
\(2n^{d/4}+O(1)\le n^{d/2}\). It would satisfy
Lemma~\ref{lem:small-grid-purity} with \(p=h/2\), a contradiction.
Stabilization therefore makes this purity property absent.

With purity absent, Lemma~\ref{lem:bias-purity} rules out availability
of absolute bias \(\ge n^{-h}\) at widths \(n^\gamma\). Stabilization
makes that property absent as well. Choosing a fixed
\(0<\eta_0\le\min(\gamma,h)\) therefore gives
\eqref{eq:source-2} for all laws on the retained sets and either color,
after a subsequence and a total of \(o(N)\) discards.
\end{proof}
